% F12 -- The group label predicts difficulty, and it can be read from the past alone.
\begin{figure}[t]
\centering
\replegend{\swatch{clmove}~average site-change rate over the day \qquad
           \swatch{clbase}~persistence baseline \qquad
           \swatch{clmodel}~detection agreement}
\plotlink{\nbllm}{\fitwidth{%
\begin{tikzpicture}
\begin{groupplot}[
  group style={group size=2 by 1, horizontal sep=2.1cm},
  repbar, width=0.42\linewidth, height=5.4cm, xtick=data,
  title style={font=\footnotesize,yshift=1pt},
]
% ---- (a) mobility vs predictability ----
\nextgroupplot[title={(a) the ranking is exactly inverted},
  symbolic x coords={G0,G1,G2,G3},
  ymin=0, ymax=78, ytick={0,20,40,60},
  ylabel={\%}, /pgf/bar width=10pt]
\addplot[fill=clmove,draw=none] coordinates {(G0,33)(G1,42)(G2,51)(G3,52)};
\addplot[fill=clbase,draw=none] coordinates {(G0,56.3)(G1,41.7)(G2,19.7)(G3,17.8)};
%
% ---- (b) detection from the context prefix ----
\nextgroupplot[title={(b) group recovered from the past alone},
  symbolic x coords={ca,cb,cc,cd},
  xticklabels={30\%,50\%,80\%,90\%},
  xlabel={context fraction at question time},
  ymin=0, ymax=118, ytick={0,25,...,100},
  ylabel={agreement (\%)}, /pgf/bar width=13pt]
\addplot[fill=clnonsig,draw=none] coordinates {(ca,44)(cb,61)};
\addplot[fill=clmodel,draw=none] coordinates {(cc,92)(cd,92)};
\node[font=\scriptsize,text=clmodel,anchor=south,align=center] at (axis cs:cc,104) {reference protocol};
\end{groupplot}
\end{tikzpicture}}}
\caption{Group difficulty and group detectability}
\label{fig:group-difficulty}
\figtext{%
(a)~The group label is not merely descriptive: it says in advance how hard a
user will be. The orange bar is a group's average site-change rate over the
day and the grey-blue bar its persistence baseline, read against each other
rather than on a common scale. G0 changes site about a third of the time and
``repeat the last cell'' is already right 56.3\,\% of the time for them; G3
changes site half again as often and the same rule is right only 17.8\,\% of
the time. The ranking of persistence is exactly the inverse of the ranking of
mobility, which is worth verifying rather than assuming. (b)~At inference only
a user's past exists, so the group must be computed from the context prefix
and never from the target, or the whole idea is circular. The bars give how
often the group detected from the prefix alone agrees with the group computed
from the full sequence, at four context fractions. At the reference 80\,\%
context the pipeline recovers the right group for 92 of the 100 held-out users
unaided, and 90\,\% context gives the same figure; below roughly 60\,\% of the
day it collapses, 61\,\% agreement at half a day and 44\,\% at 30\,\%, which
is what the two pale bars mean. The mechanism is therefore worth using on a
user with most of a day behind them and worth nothing on a cold start, which
is exactly the case in which an operator has the least information and would
most want a prediction. \smob{}, 400 train / 100 test users.}
\figsource{\nbllm}{train\_cellid\_llm.ipynb}{the group-detection and per-group baseline cells}
\end{figure}
